\documentclass[../../main.tex]{subfiles}
\begin{document}

\begin{flushright}
\footnotesize\textit{【自動文字起こし・要確認】}
\end{flushright}

{\bf [解]} $P, Q$ が2辺 $CA, AB$ 上にある時，$P$ が $CA$，$Q$ が $AB$ 上にあるとして良い．
$|PA|=x, |QA|=y \quad (0 \le x \le 36, 0 \le y \le 25 \quad \cdots \text{①})$ とする．題意から
\begin{align}
xy = \dfrac{1}{2} \cdot 25 \cdot 36 = 450 \quad \cdots \text{②}
\end{align}
$\triangle APQ$ で余弦定理から，
\begin{align}
|PQ|^2 &= |AP|^2 + |AQ|^2 - 2 AP \cdot AQ \cos \angle PAQ \\
&= x^2 + y^2 - 900 \cos \angle PAQ
\end{align}
$x^2, y^2 > 0$ から，AM-GM 及び ② から，
\begin{align}
|PQ|^2 \ge 900 (1 - \cos \angle PAQ) \quad \cdots \text{③}
\end{align}
等号成立は $x=y=15\sqrt{2}$ の時成立する．ここで $\triangle ABC$ について余弦定理から，
\begin{align}
\cos \angle CAB = \dfrac{36^2 + 25^2 - 32^2}{2 \cdot 36 \cdot 25} = \dfrac{897}{2 \cdot 900}
\end{align}
から，③より，
\begin{align}
|PQ|^2 \ge \dfrac{903}{2} \quad \cdots \text{④}
\end{align}
以下同様にして，
\begin{itemize}
  \item $P, Q$ が2辺 $AB, BC$ 上にある時
  \begin{align}
  |PQ|^2 \ge \dfrac{1247}{2} \quad (\text{等号成立は } |PB|=|QB|=20) \quad \cdots \text{⑤}
  \end{align}
  \item $P, Q$ が2辺 $BC, CA$ 上にある時
  \begin{align}
  |PQ|^2 \ge \dfrac{609}{2} \quad (\text{等号成立は } |PC|=|QC|=24) \quad \cdots \text{⑥}
  \end{align}
\end{itemize}
④，⑤，⑥から，もとめるのは，
\begin{align}
\text{2辺 } BC, CA \text{上の，} C \text{からの距離が } 24 \text{である2点．}
\end{align}

\end{document}
